\section{Removing the cap and proving the consequences}
\label{sec:uncapping}

Let $\delta$ be the uniform assignment-falsification constant from
assumption~(iv), and put
\[
    r=1-\delta<1.
\]
For every fixed assignment, independence across clause times gives
\[
    \Prb(\text{a fixed assignment satisfies }F_m)\le r^m.
\]
A union bound over the $2^n$ assignments yields
\begin{equation}
    \Prb(H_n>m)\le 2^n r^m.
    \label{eq:first-moment-tail}
\end{equation}
In particular, $H_n$ is almost surely finite and has a finite second moment.

Choose
\begin{equation}
    L_n=
    \left\lceil
      \frac{\log2+1}{\delta}\,n
    \right\rceil,
    \qquad
    T_n=H_n\wedge L_n.
    \label{eq:uncapping-level}
\end{equation}
Since $r\le e^{-\delta}$,
\[
    2^n r^{L_n}\le e^{-n}.
\]
The integer-valued tail-sum identity gives
\begin{align}
    \E(H_n-T_n)^2
    &=\sum_{j=0}^\infty(2j+1)\Prb(H_n>L_n+j)
      \notag\\
    &\le
      2^n r^{L_n}
      \sum_{j=0}^\infty(2j+1)r^j
      \notag\\
    &=2^n r^{L_n}\frac{1+r}{(1-r)^2}
      =O_{\delta}(e^{-n}).
    \label{eq:cap-error}
\end{align}
The cap in \eqref{eq:uncapping-level} is a fixed linear multiple of $n$.
Thus \eqref{eq:capped-variance} gives $\Var(T_n)=O(n)$ for
$n\ge2k$, with dependence on the local parameters.  Consequently,
\[
    \Var(H_n)
    \le
    2\Var(T_n)+2\Var(H_n-T_n)
    \le
    2\Var(T_n)+2\E(H_n-T_n)^2
    =O(n).
\]

For the finitely many sizes $k\le n<2k$, equation
\eqref{eq:first-moment-tail} directly bounds $\E H_n^2$ by a constant
depending only on $\delta$ and $k$.  Enlarging the constants proves the
uncapped assertion of Theorem~\ref{thm:local-criterion}.  The parameter
choices in \eqref{eq:local-parameters} then prove
Theorem~\ref{thm:main}.

\begin{proof}[Proof of Corollary~\ref{cor:central-window}]
Let $\mu_n=\E H_n$, and let $K$ be the constant in
Theorem~\ref{thm:main}.  Put
\[
    t_n=\sqrt{\frac{2Kn}{\varepsilon}}.
\]
Chebyshev's inequality gives
\[
    \Prb(|H_n-\mu_n|\ge t_n)\le\frac{\varepsilon}{2}.
\]
At $m=\lceil\mu_n+t_n\rceil$, this implies $s_n(m)\le\varepsilon/2$, so
$r_n(\varepsilon)\le\lceil\mu_n+t_n\rceil$.  For every nonnegative integer
$m\le\mu_n-t_n$, one has
$s_n(m)\ge1-\varepsilon/2>1-\varepsilon$, so
$r_n(1-\varepsilon)>\mu_n-t_n$ unless the latter quantity is negative, in
which case the trivial bound $r_n(1-\varepsilon)\ge0$ suffices.  Hence
\[
    r_n(\varepsilon)-r_n(1-\varepsilon)
    \le 2t_n+2.
\]
\end{proof}
